\documentclass[12pt,a4paper]{article}
\usepackage[top=1.8cm,bottom=1.8cm,left=1.8cm,right=1.8cm]{geometry}
\usepackage{fontspec}
\setmainfont{TeX Gyre Termes}
\usepackage{amsmath,amssymb}
\usepackage{fancyhdr}
\usepackage{tasks}
\usepackage{tikz}
\settasks{label=\textbf{\Alph*.},label-width=1.5em,item-indent=2em}
\setlength{\parindent}{0pt}
\setlength{\parskip}{3pt}
\pagestyle{fancy}\fancyhf{}\renewcommand{\headrulewidth}{0pt}
\fancyhead[L]{\small Ngân hàng bài tập Toán 8 -- 8H2-C7-B28-D02}
\cfoot{\small Trang \thepage}
\newif\ifloigiai
\loigiaitrue
\newenvironment{cauhoi}[2]{%
  \par\medskip\noindent\rule{\textwidth}{0.4pt}\par\nopagebreak
  \noindent\textbf{[#1]}\hfill{\scriptsize\ttfamily #2}\par\nopagebreak\smallskip}{\par}
\newcommand{\loigiai}[1]{\ifloigiai\par\smallskip\noindent\textit{Lời giải.} #1\par\fi}
\newcommand{\dapan}[1]{\ifloigiai\par\smallskip\noindent\textit{Đáp án:} \textbf{#1}\par\fi}
\begin{document}
\begin{center}
{\Large\textbf{NGÂN HÀNG BÀI TẬP TOÁN 8}}\\[3pt]
\textbf{Chương 7. Phương trình bậc nhất và hàm số bậc nhất}\\[2pt]
\textbf{Bài 28. Hàm số bậc nhất và đồ thị của hàm số bậc nhất}\\[2pt]
\textbf{Dạng D02: Vẽ đồ thị, giao điểm hai đường thẳng}\\[2pt]
\textit{Mã dạng: 8H2-C7-B28-D02 \quad 50 bài tự luận + 50 câu trắc nghiệm \quad (TH: 50)}
\end{center}
\textbf{Quy ước mã câu:} \textit{khối, học kì, chương, bài, dạng, loại (TL/TN), mức độ, số thứ tự}.\par Ví dụ: \texttt{8H2-C7-B28-D02-TL-NB-001}.\par Để in bản cho học sinh, đặt \verb|\loigiaifalse| ngay sau \verb|\loigiaitrue|.
\section*{Phần I. Tự luận}
\begin{cauhoi}{Bài 1}{8H2-C7-B28-D02-TL-TH-001}
Tìm toạ độ giao điểm của hai đường thẳng $(d_1): y=-4x-3$ và $(d_2): y=3x+11$.
\loigiai{Hoành độ giao điểm là nghiệm của phương trình $-4x-3=3x+11$, suy ra $-7x=14$ nên $x=-2$. Thay vào $(d_1)$: $y=-4\cdot(-2)-3=5$. Vậy giao điểm là $(-2;5)$.}
\end{cauhoi}
\begin{cauhoi}{Bài 2}{8H2-C7-B28-D02-TL-TH-002}
Tìm toạ độ giao điểm của hai đường thẳng $(d_1): y=4x-10$ và $(d_2): y=3x-6$.
\loigiai{Hoành độ giao điểm là nghiệm của phương trình $4x-10=3x-6$, suy ra $x=4$ nên $x=4$. Thay vào $(d_1)$: $y=4\cdot(4)-10=6$. Vậy giao điểm là $(4;6)$.}
\end{cauhoi}
\begin{cauhoi}{Bài 3}{8H2-C7-B28-D02-TL-TH-003}
Tìm toạ độ giao điểm của hai đường thẳng $(d_1): y=-2x-13$ và $(d_2): y=-4x-21$.
\loigiai{Hoành độ giao điểm là nghiệm của phương trình $-2x-13=-4x-21$, suy ra $2x=-8$ nên $x=-4$. Thay vào $(d_1)$: $y=-2\cdot(-4)-13=-5$. Vậy giao điểm là $(-4;-5)$.}
\end{cauhoi}
\begin{cauhoi}{Bài 4}{8H2-C7-B28-D02-TL-TH-004}
Tìm toạ độ giao điểm của hai đường thẳng $(d_1): y=-2x+6$ và $(d_2): y=-4x+8$.
\loigiai{Hoành độ giao điểm là nghiệm của phương trình $-2x+6=-4x+8$, suy ra $2x=2$ nên $x=1$. Thay vào $(d_1)$: $y=-2\cdot(1)+6=4$. Vậy giao điểm là $(1;4)$.}
\end{cauhoi}
\begin{cauhoi}{Bài 5}{8H2-C7-B28-D02-TL-TH-005}
Tìm toạ độ giao điểm của hai đường thẳng $(d_1): y=-2x+10$ và $(d_2): y=-4x+16$.
\loigiai{Hoành độ giao điểm là nghiệm của phương trình $-2x+10=-4x+16$, suy ra $2x=6$ nên $x=3$. Thay vào $(d_1)$: $y=-2\cdot(3)+10=4$. Vậy giao điểm là $(3;4)$.}
\end{cauhoi}
\begin{cauhoi}{Bài 6}{8H2-C7-B28-D02-TL-TH-006}
Tìm toạ độ giao điểm của hai đường thẳng $(d_1): y=-3x-8$ và $(d_2): y=2x+17$.
\loigiai{Hoành độ giao điểm là nghiệm của phương trình $-3x-8=2x+17$, suy ra $-5x=25$ nên $x=-5$. Thay vào $(d_1)$: $y=-3\cdot(-5)-8=7$. Vậy giao điểm là $(-5;7)$.}
\end{cauhoi}
\begin{cauhoi}{Bài 7}{8H2-C7-B28-D02-TL-TH-007}
Tìm toạ độ giao điểm của hai đường thẳng $(d_1): y=-4x+21$ và $(d_2): y=-2x+9$.
\loigiai{Hoành độ giao điểm là nghiệm của phương trình $-4x+21=-2x+9$, suy ra $-2x=-12$ nên $x=6$. Thay vào $(d_1)$: $y=-4\cdot(6)+21=-3$. Vậy giao điểm là $(6;-3)$.}
\end{cauhoi}
\begin{cauhoi}{Bài 8}{8H2-C7-B28-D02-TL-TH-008}
Tìm toạ độ giao điểm của hai đường thẳng $(d_1): y=-3x-3$ và $(d_2): y=x+5$.
\loigiai{Hoành độ giao điểm là nghiệm của phương trình $-3x-3=x+5$, suy ra $-4x=8$ nên $x=-2$. Thay vào $(d_1)$: $y=-3\cdot(-2)-3=3$. Vậy giao điểm là $(-2;3)$.}
\end{cauhoi}
\begin{cauhoi}{Bài 9}{8H2-C7-B28-D02-TL-TH-009}
Tìm toạ độ giao điểm của hai đường thẳng $(d_1): y=-4x+26$ và $(d_2): y=2x-4$.
\loigiai{Hoành độ giao điểm là nghiệm của phương trình $-4x+26=2x-4$, suy ra $-6x=-30$ nên $x=5$. Thay vào $(d_1)$: $y=-4\cdot(5)+26=6$. Vậy giao điểm là $(5;6)$.}
\end{cauhoi}
\begin{cauhoi}{Bài 10}{8H2-C7-B28-D02-TL-TH-010}
Tìm toạ độ giao điểm của hai đường thẳng $(d_1): y=-2x+4$ và $(d_2): y=-4x+2$.
\loigiai{Hoành độ giao điểm là nghiệm của phương trình $-2x+4=-4x+2$, suy ra $2x=-2$ nên $x=-1$. Thay vào $(d_1)$: $y=-2\cdot(-1)+4=6$. Vậy giao điểm là $(-1;6)$.}
\end{cauhoi}
\begin{cauhoi}{Bài 11}{8H2-C7-B28-D02-TL-TH-011}
Tìm toạ độ giao điểm của hai đường thẳng $(d_1): y=-x-6$ và $(d_2): y=-4x-15$.
\loigiai{Hoành độ giao điểm là nghiệm của phương trình $-x-6=-4x-15$, suy ra $3x=-9$ nên $x=-3$. Thay vào $(d_1)$: $y=-1\cdot(-3)-6=-3$. Vậy giao điểm là $(-3;-3)$.}
\end{cauhoi}
\begin{cauhoi}{Bài 12}{8H2-C7-B28-D02-TL-TH-012}
Tìm toạ độ giao điểm của hai đường thẳng $(d_1): y=4x-9$ và $(d_2): y=-3x+5$.
\loigiai{Hoành độ giao điểm là nghiệm của phương trình $4x-9=-3x+5$, suy ra $7x=14$ nên $x=2$. Thay vào $(d_1)$: $y=4\cdot(2)-9=-1$. Vậy giao điểm là $(2;-1)$.}
\end{cauhoi}
\begin{cauhoi}{Bài 13}{8H2-C7-B28-D02-TL-TH-013}
Tìm toạ độ giao điểm của hai đường thẳng $(d_1): y=-2x+1$ và $(d_2): y=4x+19$.
\loigiai{Hoành độ giao điểm là nghiệm của phương trình $-2x+1=4x+19$, suy ra $-6x=18$ nên $x=-3$. Thay vào $(d_1)$: $y=-2\cdot(-3)+1=7$. Vậy giao điểm là $(-3;7)$.}
\end{cauhoi}
\begin{cauhoi}{Bài 14}{8H2-C7-B28-D02-TL-TH-014}
Tìm toạ độ giao điểm của hai đường thẳng $(d_1): y=x+1$ và $(d_2): y=-x-5$.
\loigiai{Hoành độ giao điểm là nghiệm của phương trình $x+1=-x-5$, suy ra $2x=-6$ nên $x=-3$. Thay vào $(d_1)$: $y=1\cdot(-3)+1=-2$. Vậy giao điểm là $(-3;-2)$.}
\end{cauhoi}
\begin{cauhoi}{Bài 15}{8H2-C7-B28-D02-TL-TH-015}
Tìm toạ độ giao điểm của hai đường thẳng $(d_1): y=-3x+1$ và $(d_2): y=-4x+1$.
\loigiai{Hoành độ giao điểm là nghiệm của phương trình $-3x+1=-4x+1$, suy ra $x=0$ nên $x=0$. Thay vào $(d_1)$: $y=-3\cdot(0)+1=1$. Vậy giao điểm là $(0;1)$.}
\end{cauhoi}
\begin{cauhoi}{Bài 16}{8H2-C7-B28-D02-TL-TH-016}
Tìm toạ độ giao điểm của hai đường thẳng $(d_1): y=-3x-9$ và $(d_2): y=2x+11$.
\loigiai{Hoành độ giao điểm là nghiệm của phương trình $-3x-9=2x+11$, suy ra $-5x=20$ nên $x=-4$. Thay vào $(d_1)$: $y=-3\cdot(-4)-9=3$. Vậy giao điểm là $(-4;3)$.}
\end{cauhoi}
\begin{cauhoi}{Bài 17}{8H2-C7-B28-D02-TL-TH-017}
Tìm toạ độ giao điểm của hai đường thẳng $(d_1): y=-3x-9$ và $(d_2): y=-x-3$.
\loigiai{Hoành độ giao điểm là nghiệm của phương trình $-3x-9=-x-3$, suy ra $-2x=6$ nên $x=-3$. Thay vào $(d_1)$: $y=-3\cdot(-3)-9=0$. Vậy giao điểm là $(-3;0)$.}
\end{cauhoi}
\begin{cauhoi}{Bài 18}{8H2-C7-B28-D02-TL-TH-018}
Tìm toạ độ giao điểm của hai đường thẳng $(d_1): y=-2x+16$ và $(d_2): y=4x-8$.
\loigiai{Hoành độ giao điểm là nghiệm của phương trình $-2x+16=4x-8$, suy ra $-6x=-24$ nên $x=4$. Thay vào $(d_1)$: $y=-2\cdot(4)+16=8$. Vậy giao điểm là $(4;8)$.}
\end{cauhoi}
\begin{cauhoi}{Bài 19}{8H2-C7-B28-D02-TL-TH-019}
Tìm toạ độ giao điểm của hai đường thẳng $(d_1): y=2x-9$ và $(d_2): y=-x+6$.
\loigiai{Hoành độ giao điểm là nghiệm của phương trình $2x-9=-x+6$, suy ra $3x=15$ nên $x=5$. Thay vào $(d_1)$: $y=2\cdot(5)-9=1$. Vậy giao điểm là $(5;1)$.}
\end{cauhoi}
\begin{cauhoi}{Bài 20}{8H2-C7-B28-D02-TL-TH-020}
Tìm toạ độ giao điểm của hai đường thẳng $(d_1): y=-4x+20$ và $(d_2): y=x$.
\loigiai{Hoành độ giao điểm là nghiệm của phương trình $-4x+20=x$, suy ra $-5x=-20$ nên $x=4$. Thay vào $(d_1)$: $y=-4\cdot(4)+20=4$. Vậy giao điểm là $(4;4)$.}
\end{cauhoi}
\begin{cauhoi}{Bài 21}{8H2-C7-B28-D02-TL-TH-021}
Tìm toạ độ giao điểm của hai đường thẳng $(d_1): y=-4x+18$ và $(d_2): y=2x-18$.
\loigiai{Hoành độ giao điểm là nghiệm của phương trình $-4x+18=2x-18$, suy ra $-6x=-36$ nên $x=6$. Thay vào $(d_1)$: $y=-4\cdot(6)+18=-6$. Vậy giao điểm là $(6;-6)$.}
\end{cauhoi}
\begin{cauhoi}{Bài 22}{8H2-C7-B28-D02-TL-TH-022}
Tìm toạ độ giao điểm của hai đường thẳng $(d_1): y=-3x+21$ và $(d_2): y=3x-9$.
\loigiai{Hoành độ giao điểm là nghiệm của phương trình $-3x+21=3x-9$, suy ra $-6x=-30$ nên $x=5$. Thay vào $(d_1)$: $y=-3\cdot(5)+21=6$. Vậy giao điểm là $(5;6)$.}
\end{cauhoi}
\begin{cauhoi}{Bài 23}{8H2-C7-B28-D02-TL-TH-023}
Tìm toạ độ giao điểm của hai đường thẳng $(d_1): y=-x$ và $(d_2): y=-2x+1$.
\loigiai{Hoành độ giao điểm là nghiệm của phương trình $-x=-2x+1$, suy ra $x=1$ nên $x=1$. Thay vào $(d_1)$: $y=-1\cdot(1)+0=-1$. Vậy giao điểm là $(1;-1)$.}
\end{cauhoi}
\begin{cauhoi}{Bài 24}{8H2-C7-B28-D02-TL-TH-024}
Tìm toạ độ giao điểm của hai đường thẳng $(d_1): y=x$ và $(d_2): y=2x+5$.
\loigiai{Hoành độ giao điểm là nghiệm của phương trình $x=2x+5$, suy ra $-x=5$ nên $x=-5$. Thay vào $(d_1)$: $y=1\cdot(-5)+0=-5$. Vậy giao điểm là $(-5;-5)$.}
\end{cauhoi}
\begin{cauhoi}{Bài 25}{8H2-C7-B28-D02-TL-TH-025}
Tìm toạ độ giao điểm của hai đường thẳng $(d_1): y=-3x-7$ và $(d_2): y=3x+17$.
\loigiai{Hoành độ giao điểm là nghiệm của phương trình $-3x-7=3x+17$, suy ra $-6x=24$ nên $x=-4$. Thay vào $(d_1)$: $y=-3\cdot(-4)-7=5$. Vậy giao điểm là $(-4;5)$.}
\end{cauhoi}
\begin{cauhoi}{Bài 26}{8H2-C7-B28-D02-TL-TH-026}
Tìm toạ độ giao điểm của hai đường thẳng $(d_1): y=-4x-2$ và $(d_2): y=2x-2$.
\loigiai{Hoành độ giao điểm là nghiệm của phương trình $-4x-2=2x-2$, suy ra $-6x=0$ nên $x=0$. Thay vào $(d_1)$: $y=-4\cdot(0)-2=-2$. Vậy giao điểm là $(0;-2)$.}
\end{cauhoi}
\begin{cauhoi}{Bài 27}{8H2-C7-B28-D02-TL-TH-027}
Tìm toạ độ giao điểm của hai đường thẳng $(d_1): y=-3x-13$ và $(d_2): y=-2x-8$.
\loigiai{Hoành độ giao điểm là nghiệm của phương trình $-3x-13=-2x-8$, suy ra $-x=5$ nên $x=-5$. Thay vào $(d_1)$: $y=-3\cdot(-5)-13=2$. Vậy giao điểm là $(-5;2)$.}
\end{cauhoi}
\begin{cauhoi}{Bài 28}{8H2-C7-B28-D02-TL-TH-028}
Tìm toạ độ giao điểm của hai đường thẳng $(d_1): y=3x-8$ và $(d_2): y=-3x+10$.
\loigiai{Hoành độ giao điểm là nghiệm của phương trình $3x-8=-3x+10$, suy ra $6x=18$ nên $x=3$. Thay vào $(d_1)$: $y=3\cdot(3)-8=1$. Vậy giao điểm là $(3;1)$.}
\end{cauhoi}
\begin{cauhoi}{Bài 29}{8H2-C7-B28-D02-TL-TH-029}
Tìm toạ độ giao điểm của hai đường thẳng $(d_1): y=2x+6$ và $(d_2): y=-2x-2$.
\loigiai{Hoành độ giao điểm là nghiệm của phương trình $2x+6=-2x-2$, suy ra $4x=-8$ nên $x=-2$. Thay vào $(d_1)$: $y=2\cdot(-2)+6=2$. Vậy giao điểm là $(-2;2)$.}
\end{cauhoi}
\begin{cauhoi}{Bài 30}{8H2-C7-B28-D02-TL-TH-030}
Tìm toạ độ giao điểm của hai đường thẳng $(d_1): y=2x+14$ và $(d_2): y=x+9$.
\loigiai{Hoành độ giao điểm là nghiệm của phương trình $2x+14=x+9$, suy ra $x=-5$ nên $x=-5$. Thay vào $(d_1)$: $y=2\cdot(-5)+14=4$. Vậy giao điểm là $(-5;4)$.}
\end{cauhoi}
\begin{cauhoi}{Bài 31}{8H2-C7-B28-D02-TL-TH-031}
Tìm toạ độ giao điểm của hai đường thẳng $(d_1): y=-x+4$ và $(d_2): y=x+6$.
\loigiai{Hoành độ giao điểm là nghiệm của phương trình $-x+4=x+6$, suy ra $-2x=2$ nên $x=-1$. Thay vào $(d_1)$: $y=-1\cdot(-1)+4=5$. Vậy giao điểm là $(-1;5)$.}
\end{cauhoi}
\begin{cauhoi}{Bài 32}{8H2-C7-B28-D02-TL-TH-032}
Tìm toạ độ giao điểm của hai đường thẳng $(d_1): y=x+9$ và $(d_2): y=-3x-3$.
\loigiai{Hoành độ giao điểm là nghiệm của phương trình $x+9=-3x-3$, suy ra $4x=-12$ nên $x=-3$. Thay vào $(d_1)$: $y=1\cdot(-3)+9=6$. Vậy giao điểm là $(-3;6)$.}
\end{cauhoi}
\begin{cauhoi}{Bài 33}{8H2-C7-B28-D02-TL-TH-033}
Tìm toạ độ giao điểm của hai đường thẳng $(d_1): y=-2x+2$ và $(d_2): y=-3x+5$.
\loigiai{Hoành độ giao điểm là nghiệm của phương trình $-2x+2=-3x+5$, suy ra $x=3$ nên $x=3$. Thay vào $(d_1)$: $y=-2\cdot(3)+2=-4$. Vậy giao điểm là $(3;-4)$.}
\end{cauhoi}
\begin{cauhoi}{Bài 34}{8H2-C7-B28-D02-TL-TH-034}
Tìm toạ độ giao điểm của hai đường thẳng $(d_1): y=3x-1$ và $(d_2): y=-4x+6$.
\loigiai{Hoành độ giao điểm là nghiệm của phương trình $3x-1=-4x+6$, suy ra $7x=7$ nên $x=1$. Thay vào $(d_1)$: $y=3\cdot(1)-1=2$. Vậy giao điểm là $(1;2)$.}
\end{cauhoi}
\begin{cauhoi}{Bài 35}{8H2-C7-B28-D02-TL-TH-035}
Tìm toạ độ giao điểm của hai đường thẳng $(d_1): y=-x+7$ và $(d_2): y=4x+12$.
\loigiai{Hoành độ giao điểm là nghiệm của phương trình $-x+7=4x+12$, suy ra $-5x=5$ nên $x=-1$. Thay vào $(d_1)$: $y=-1\cdot(-1)+7=8$. Vậy giao điểm là $(-1;8)$.}
\end{cauhoi}
\begin{cauhoi}{Bài 36}{8H2-C7-B28-D02-TL-TH-036}
Tìm toạ độ giao điểm của hai đường thẳng $(d_1): y=4x+6$ và $(d_2): y=-4x+6$.
\loigiai{Hoành độ giao điểm là nghiệm của phương trình $4x+6=-4x+6$, suy ra $8x=0$ nên $x=0$. Thay vào $(d_1)$: $y=4\cdot(0)+6=6$. Vậy giao điểm là $(0;6)$.}
\end{cauhoi}
\begin{cauhoi}{Bài 37}{8H2-C7-B28-D02-TL-TH-037}
Tìm toạ độ giao điểm của hai đường thẳng $(d_1): y=-2x+8$ và $(d_2): y=-4x+10$.
\loigiai{Hoành độ giao điểm là nghiệm của phương trình $-2x+8=-4x+10$, suy ra $2x=2$ nên $x=1$. Thay vào $(d_1)$: $y=-2\cdot(1)+8=6$. Vậy giao điểm là $(1;6)$.}
\end{cauhoi}
\begin{cauhoi}{Bài 38}{8H2-C7-B28-D02-TL-TH-038}
Tìm toạ độ giao điểm của hai đường thẳng $(d_1): y=3x-4$ và $(d_2): y=-3x+2$.
\loigiai{Hoành độ giao điểm là nghiệm của phương trình $3x-4=-3x+2$, suy ra $6x=6$ nên $x=1$. Thay vào $(d_1)$: $y=3\cdot(1)-4=-1$. Vậy giao điểm là $(1;-1)$.}
\end{cauhoi}
\begin{cauhoi}{Bài 39}{8H2-C7-B28-D02-TL-TH-039}
Tìm toạ độ giao điểm của hai đường thẳng $(d_1): y=4x+24$ và $(d_2): y=2x+14$.
\loigiai{Hoành độ giao điểm là nghiệm của phương trình $4x+24=2x+14$, suy ra $2x=-10$ nên $x=-5$. Thay vào $(d_1)$: $y=4\cdot(-5)+24=4$. Vậy giao điểm là $(-5;4)$.}
\end{cauhoi}
\begin{cauhoi}{Bài 40}{8H2-C7-B28-D02-TL-TH-040}
Tìm toạ độ giao điểm của hai đường thẳng $(d_1): y=-2x-3$ và $(d_2): y=x+9$.
\loigiai{Hoành độ giao điểm là nghiệm của phương trình $-2x-3=x+9$, suy ra $-3x=12$ nên $x=-4$. Thay vào $(d_1)$: $y=-2\cdot(-4)-3=5$. Vậy giao điểm là $(-4;5)$.}
\end{cauhoi}
\begin{cauhoi}{Bài 41}{8H2-C7-B28-D02-TL-TH-041}
Tìm toạ độ giao điểm của hai đường thẳng $(d_1): y=-x+4$ và $(d_2): y=2x+4$.
\loigiai{Hoành độ giao điểm là nghiệm của phương trình $-x+4=2x+4$, suy ra $-3x=0$ nên $x=0$. Thay vào $(d_1)$: $y=-1\cdot(0)+4=4$. Vậy giao điểm là $(0;4)$.}
\end{cauhoi}
\begin{cauhoi}{Bài 42}{8H2-C7-B28-D02-TL-TH-042}
Tìm toạ độ giao điểm của hai đường thẳng $(d_1): y=-x+7$ và $(d_2): y=2x+10$.
\loigiai{Hoành độ giao điểm là nghiệm của phương trình $-x+7=2x+10$, suy ra $-3x=3$ nên $x=-1$. Thay vào $(d_1)$: $y=-1\cdot(-1)+7=8$. Vậy giao điểm là $(-1;8)$.}
\end{cauhoi}
\begin{cauhoi}{Bài 43}{8H2-C7-B28-D02-TL-TH-043}
Tìm toạ độ giao điểm của hai đường thẳng $(d_1): y=4x+8$ và $(d_2): y=-x-2$.
\loigiai{Hoành độ giao điểm là nghiệm của phương trình $4x+8=-x-2$, suy ra $5x=-10$ nên $x=-2$. Thay vào $(d_1)$: $y=4\cdot(-2)+8=0$. Vậy giao điểm là $(-2;0)$.}
\end{cauhoi}
\begin{cauhoi}{Bài 44}{8H2-C7-B28-D02-TL-TH-044}
Tìm toạ độ giao điểm của hai đường thẳng $(d_1): y=-x-4$ và $(d_2): y=2x+2$.
\loigiai{Hoành độ giao điểm là nghiệm của phương trình $-x-4=2x+2$, suy ra $-3x=6$ nên $x=-2$. Thay vào $(d_1)$: $y=-1\cdot(-2)-4=-2$. Vậy giao điểm là $(-2;-2)$.}
\end{cauhoi}
\begin{cauhoi}{Bài 45}{8H2-C7-B28-D02-TL-TH-045}
Tìm toạ độ giao điểm của hai đường thẳng $(d_1): y=4x+20$ và $(d_2): y=-4x-20$.
\loigiai{Hoành độ giao điểm là nghiệm của phương trình $4x+20=-4x-20$, suy ra $8x=-40$ nên $x=-5$. Thay vào $(d_1)$: $y=4\cdot(-5)+20=0$. Vậy giao điểm là $(-5;0)$.}
\end{cauhoi}
\begin{cauhoi}{Bài 46}{8H2-C7-B28-D02-TL-TH-046}
Tìm toạ độ giao điểm của hai đường thẳng $(d_1): y=x-2$ và $(d_2): y=-3x+6$.
\loigiai{Hoành độ giao điểm là nghiệm của phương trình $x-2=-3x+6$, suy ra $4x=8$ nên $x=2$. Thay vào $(d_1)$: $y=1\cdot(2)-2=0$. Vậy giao điểm là $(2;0)$.}
\end{cauhoi}
\begin{cauhoi}{Bài 47}{8H2-C7-B28-D02-TL-TH-047}
Tìm toạ độ giao điểm của hai đường thẳng $(d_1): y=-2x+12$ và $(d_2): y=-3x+16$.
\loigiai{Hoành độ giao điểm là nghiệm của phương trình $-2x+12=-3x+16$, suy ra $x=4$ nên $x=4$. Thay vào $(d_1)$: $y=-2\cdot(4)+12=4$. Vậy giao điểm là $(4;4)$.}
\end{cauhoi}
\begin{cauhoi}{Bài 48}{8H2-C7-B28-D02-TL-TH-048}
Tìm toạ độ giao điểm của hai đường thẳng $(d_1): y=3x+14$ và $(d_2): y=2x+10$.
\loigiai{Hoành độ giao điểm là nghiệm của phương trình $3x+14=2x+10$, suy ra $x=-4$ nên $x=-4$. Thay vào $(d_1)$: $y=3\cdot(-4)+14=2$. Vậy giao điểm là $(-4;2)$.}
\end{cauhoi}
\begin{cauhoi}{Bài 49}{8H2-C7-B28-D02-TL-TH-049}
Tìm toạ độ giao điểm của hai đường thẳng $(d_1): y=x+8$ và $(d_2): y=2x+8$.
\loigiai{Hoành độ giao điểm là nghiệm của phương trình $x+8=2x+8$, suy ra $-x=0$ nên $x=0$. Thay vào $(d_1)$: $y=1\cdot(0)+8=8$. Vậy giao điểm là $(0;8)$.}
\end{cauhoi}
\begin{cauhoi}{Bài 50}{8H2-C7-B28-D02-TL-TH-050}
Tìm toạ độ giao điểm của hai đường thẳng $(d_1): y=3x-13$ và $(d_2): y=x-5$.
\loigiai{Hoành độ giao điểm là nghiệm của phương trình $3x-13=x-5$, suy ra $2x=8$ nên $x=4$. Thay vào $(d_1)$: $y=3\cdot(4)-13=-1$. Vậy giao điểm là $(4;-1)$.}
\end{cauhoi}
\clearpage\section*{Phần II. Trắc nghiệm}
\begin{cauhoi}{Câu 1}{8H2-C7-B28-D02-TN-TH-001}
Tìm toạ độ giao điểm của hai đường thẳng $(d_1): y=-4x-3$ và $(d_2): y=3x+11$.
\begin{tasks}(2)
\task $(5;-2)$
\task $(2;5)$
\task $(-2;5)$
\task $(-2;-5)$
\end{tasks}
\dapan{C}
\end{cauhoi}
\begin{cauhoi}{Câu 2}{8H2-C7-B28-D02-TN-TH-002}
Tìm toạ độ giao điểm của hai đường thẳng $(d_1): y=4x-10$ và $(d_2): y=3x-6$.
\begin{tasks}(2)
\task $(4;-6)$
\task $(4;6)$
\task $(6;4)$
\task $(-4;6)$
\end{tasks}
\dapan{B}
\end{cauhoi}
\begin{cauhoi}{Câu 3}{8H2-C7-B28-D02-TN-TH-003}
Tìm toạ độ giao điểm của hai đường thẳng $(d_1): y=-2x-13$ và $(d_2): y=-4x-21$.
\begin{tasks}(2)
\task $(-4;-5)$
\task $(4;-5)$
\task $(-5;-4)$
\task $(-4;5)$
\end{tasks}
\dapan{A}
\end{cauhoi}
\begin{cauhoi}{Câu 4}{8H2-C7-B28-D02-TN-TH-004}
Tìm toạ độ giao điểm của hai đường thẳng $(d_1): y=-2x+6$ và $(d_2): y=-4x+8$.
\begin{tasks}(2)
\task $(1;4)$
\task $(1;-4)$
\task $(4;1)$
\task $(-1;4)$
\end{tasks}
\dapan{A}
\end{cauhoi}
\begin{cauhoi}{Câu 5}{8H2-C7-B28-D02-TN-TH-005}
Tìm toạ độ giao điểm của hai đường thẳng $(d_1): y=-2x+10$ và $(d_2): y=-4x+16$.
\begin{tasks}(2)
\task $(-3;4)$
\task $(3;4)$
\task $(4;3)$
\task $(3;-4)$
\end{tasks}
\dapan{B}
\end{cauhoi}
\begin{cauhoi}{Câu 6}{8H2-C7-B28-D02-TN-TH-006}
Tìm toạ độ giao điểm của hai đường thẳng $(d_1): y=-3x-8$ và $(d_2): y=2x+17$.
\begin{tasks}(2)
\task $(7;-5)$
\task $(-5;7)$
\task $(-5;-7)$
\task $(5;7)$
\end{tasks}
\dapan{B}
\end{cauhoi}
\begin{cauhoi}{Câu 7}{8H2-C7-B28-D02-TN-TH-007}
Tìm toạ độ giao điểm của hai đường thẳng $(d_1): y=-4x+21$ và $(d_2): y=-2x+9$.
\begin{tasks}(2)
\task $(-3;6)$
\task $(-6;-3)$
\task $(6;3)$
\task $(6;-3)$
\end{tasks}
\dapan{D}
\end{cauhoi}
\begin{cauhoi}{Câu 8}{8H2-C7-B28-D02-TN-TH-008}
Tìm toạ độ giao điểm của hai đường thẳng $(d_1): y=-3x-3$ và $(d_2): y=x+5$.
\begin{tasks}(2)
\task $(-2;-3)$
\task $(-2;3)$
\task $(3;-2)$
\task $(2;3)$
\end{tasks}
\dapan{B}
\end{cauhoi}
\begin{cauhoi}{Câu 9}{8H2-C7-B28-D02-TN-TH-009}
Tìm toạ độ giao điểm của hai đường thẳng $(d_1): y=-4x+26$ và $(d_2): y=2x-4$.
\begin{tasks}(2)
\task $(5;-6)$
\task $(6;5)$
\task $(5;6)$
\task $(-5;6)$
\end{tasks}
\dapan{C}
\end{cauhoi}
\begin{cauhoi}{Câu 10}{8H2-C7-B28-D02-TN-TH-010}
Tìm toạ độ giao điểm của hai đường thẳng $(d_1): y=-2x+4$ và $(d_2): y=-4x+2$.
\begin{tasks}(2)
\task $(1;6)$
\task $(-1;6)$
\task $(6;-1)$
\task $(-1;-6)$
\end{tasks}
\dapan{B}
\end{cauhoi}
\begin{cauhoi}{Câu 11}{8H2-C7-B28-D02-TN-TH-011}
Tìm toạ độ giao điểm của hai đường thẳng $(d_1): y=-x-6$ và $(d_2): y=-4x-15$.
\begin{tasks}(2)
\task $(3;-3)$
\task $(-2;-3)$
\task $(-3;-3)$
\task $(-3;3)$
\end{tasks}
\dapan{C}
\end{cauhoi}
\begin{cauhoi}{Câu 12}{8H2-C7-B28-D02-TN-TH-012}
Tìm toạ độ giao điểm của hai đường thẳng $(d_1): y=4x-9$ và $(d_2): y=-3x+5$.
\begin{tasks}(2)
\task $(2;-1)$
\task $(-1;2)$
\task $(-2;-1)$
\task $(2;1)$
\end{tasks}
\dapan{A}
\end{cauhoi}
\begin{cauhoi}{Câu 13}{8H2-C7-B28-D02-TN-TH-013}
Tìm toạ độ giao điểm của hai đường thẳng $(d_1): y=-2x+1$ và $(d_2): y=4x+19$.
\begin{tasks}(2)
\task $(-3;-7)$
\task $(-3;7)$
\task $(3;7)$
\task $(7;-3)$
\end{tasks}
\dapan{B}
\end{cauhoi}
\begin{cauhoi}{Câu 14}{8H2-C7-B28-D02-TN-TH-014}
Tìm toạ độ giao điểm của hai đường thẳng $(d_1): y=x+1$ và $(d_2): y=-x-5$.
\begin{tasks}(2)
\task $(-3;-2)$
\task $(-2;-3)$
\task $(-3;2)$
\task $(3;-2)$
\end{tasks}
\dapan{A}
\end{cauhoi}
\begin{cauhoi}{Câu 15}{8H2-C7-B28-D02-TN-TH-015}
Tìm toạ độ giao điểm của hai đường thẳng $(d_1): y=-3x+1$ và $(d_2): y=-4x+1$.
\begin{tasks}(2)
\task $(0;1)$
\task $(0;-1)$
\task $(1;1)$
\task $(1;0)$
\end{tasks}
\dapan{A}
\end{cauhoi}
\begin{cauhoi}{Câu 16}{8H2-C7-B28-D02-TN-TH-016}
Tìm toạ độ giao điểm của hai đường thẳng $(d_1): y=-3x-9$ và $(d_2): y=2x+11$.
\begin{tasks}(2)
\task $(4;3)$
\task $(-4;-3)$
\task $(-4;3)$
\task $(3;-4)$
\end{tasks}
\dapan{C}
\end{cauhoi}
\begin{cauhoi}{Câu 17}{8H2-C7-B28-D02-TN-TH-017}
Tìm toạ độ giao điểm của hai đường thẳng $(d_1): y=-3x-9$ và $(d_2): y=-x-3$.
\begin{tasks}(2)
\task $(-2;0)$
\task $(-3;0)$
\task $(3;0)$
\task $(0;-3)$
\end{tasks}
\dapan{B}
\end{cauhoi}
\begin{cauhoi}{Câu 18}{8H2-C7-B28-D02-TN-TH-018}
Tìm toạ độ giao điểm của hai đường thẳng $(d_1): y=-2x+16$ và $(d_2): y=4x-8$.
\begin{tasks}(2)
\task $(8;4)$
\task $(-4;8)$
\task $(4;8)$
\task $(4;-8)$
\end{tasks}
\dapan{C}
\end{cauhoi}
\begin{cauhoi}{Câu 19}{8H2-C7-B28-D02-TN-TH-019}
Tìm toạ độ giao điểm của hai đường thẳng $(d_1): y=2x-9$ và $(d_2): y=-x+6$.
\begin{tasks}(2)
\task $(5;1)$
\task $(5;-1)$
\task $(-5;1)$
\task $(1;5)$
\end{tasks}
\dapan{A}
\end{cauhoi}
\begin{cauhoi}{Câu 20}{8H2-C7-B28-D02-TN-TH-020}
Tìm toạ độ giao điểm của hai đường thẳng $(d_1): y=-4x+20$ và $(d_2): y=x$.
\begin{tasks}(2)
\task $(-4;4)$
\task $(5;4)$
\task $(4;-4)$
\task $(4;4)$
\end{tasks}
\dapan{D}
\end{cauhoi}
\begin{cauhoi}{Câu 21}{8H2-C7-B28-D02-TN-TH-021}
Tìm toạ độ giao điểm của hai đường thẳng $(d_1): y=-4x+18$ và $(d_2): y=2x-18$.
\begin{tasks}(2)
\task $(6;6)$
\task $(-6;6)$
\task $(6;-6)$
\task $(-6;-6)$
\end{tasks}
\dapan{C}
\end{cauhoi}
\begin{cauhoi}{Câu 22}{8H2-C7-B28-D02-TN-TH-022}
Tìm toạ độ giao điểm của hai đường thẳng $(d_1): y=-3x+21$ và $(d_2): y=3x-9$.
\begin{tasks}(2)
\task $(5;-6)$
\task $(-5;6)$
\task $(5;6)$
\task $(6;5)$
\end{tasks}
\dapan{C}
\end{cauhoi}
\begin{cauhoi}{Câu 23}{8H2-C7-B28-D02-TN-TH-023}
Tìm toạ độ giao điểm của hai đường thẳng $(d_1): y=-x$ và $(d_2): y=-2x+1$.
\begin{tasks}(2)
\task $(-1;1)$
\task $(1;1)$
\task $(1;-1)$
\task $(-1;-1)$
\end{tasks}
\dapan{C}
\end{cauhoi}
\begin{cauhoi}{Câu 24}{8H2-C7-B28-D02-TN-TH-024}
Tìm toạ độ giao điểm của hai đường thẳng $(d_1): y=x$ và $(d_2): y=2x+5$.
\begin{tasks}(2)
\task $(-4;-5)$
\task $(-5;5)$
\task $(5;-5)$
\task $(-5;-5)$
\end{tasks}
\dapan{D}
\end{cauhoi}
\begin{cauhoi}{Câu 25}{8H2-C7-B28-D02-TN-TH-025}
Tìm toạ độ giao điểm của hai đường thẳng $(d_1): y=-3x-7$ và $(d_2): y=3x+17$.
\begin{tasks}(2)
\task $(-4;-5)$
\task $(5;-4)$
\task $(4;5)$
\task $(-4;5)$
\end{tasks}
\dapan{D}
\end{cauhoi}
\begin{cauhoi}{Câu 26}{8H2-C7-B28-D02-TN-TH-026}
Tìm toạ độ giao điểm của hai đường thẳng $(d_1): y=-4x-2$ và $(d_2): y=2x-2$.
\begin{tasks}(2)
\task $(0;-2)$
\task $(-2;0)$
\task $(0;2)$
\task $(1;-2)$
\end{tasks}
\dapan{A}
\end{cauhoi}
\begin{cauhoi}{Câu 27}{8H2-C7-B28-D02-TN-TH-027}
Tìm toạ độ giao điểm của hai đường thẳng $(d_1): y=-3x-13$ và $(d_2): y=-2x-8$.
\begin{tasks}(2)
\task $(-5;-2)$
\task $(-5;2)$
\task $(5;2)$
\task $(2;-5)$
\end{tasks}
\dapan{B}
\end{cauhoi}
\begin{cauhoi}{Câu 28}{8H2-C7-B28-D02-TN-TH-028}
Tìm toạ độ giao điểm của hai đường thẳng $(d_1): y=3x-8$ và $(d_2): y=-3x+10$.
\begin{tasks}(2)
\task $(1;3)$
\task $(-3;1)$
\task $(3;1)$
\task $(3;-1)$
\end{tasks}
\dapan{C}
\end{cauhoi}
\begin{cauhoi}{Câu 29}{8H2-C7-B28-D02-TN-TH-029}
Tìm toạ độ giao điểm của hai đường thẳng $(d_1): y=2x+6$ và $(d_2): y=-2x-2$.
\begin{tasks}(2)
\task $(2;-2)$
\task $(-2;2)$
\task $(-2;-2)$
\task $(2;2)$
\end{tasks}
\dapan{B}
\end{cauhoi}
\begin{cauhoi}{Câu 30}{8H2-C7-B28-D02-TN-TH-030}
Tìm toạ độ giao điểm của hai đường thẳng $(d_1): y=2x+14$ và $(d_2): y=x+9$.
\begin{tasks}(2)
\task $(-5;4)$
\task $(-5;-4)$
\task $(5;4)$
\task $(4;-5)$
\end{tasks}
\dapan{A}
\end{cauhoi}
\begin{cauhoi}{Câu 31}{8H2-C7-B28-D02-TN-TH-031}
Tìm toạ độ giao điểm của hai đường thẳng $(d_1): y=-x+4$ và $(d_2): y=x+6$.
\begin{tasks}(2)
\task $(-1;5)$
\task $(5;-1)$
\task $(-1;-5)$
\task $(1;5)$
\end{tasks}
\dapan{A}
\end{cauhoi}
\begin{cauhoi}{Câu 32}{8H2-C7-B28-D02-TN-TH-032}
Tìm toạ độ giao điểm của hai đường thẳng $(d_1): y=x+9$ và $(d_2): y=-3x-3$.
\begin{tasks}(2)
\task $(-3;6)$
\task $(6;-3)$
\task $(-3;-6)$
\task $(3;6)$
\end{tasks}
\dapan{A}
\end{cauhoi}
\begin{cauhoi}{Câu 33}{8H2-C7-B28-D02-TN-TH-033}
Tìm toạ độ giao điểm của hai đường thẳng $(d_1): y=-2x+2$ và $(d_2): y=-3x+5$.
\begin{tasks}(2)
\task $(-4;3)$
\task $(-3;-4)$
\task $(3;-4)$
\task $(3;4)$
\end{tasks}
\dapan{C}
\end{cauhoi}
\begin{cauhoi}{Câu 34}{8H2-C7-B28-D02-TN-TH-034}
Tìm toạ độ giao điểm của hai đường thẳng $(d_1): y=3x-1$ và $(d_2): y=-4x+6$.
\begin{tasks}(2)
\task $(2;1)$
\task $(-1;2)$
\task $(1;2)$
\task $(1;-2)$
\end{tasks}
\dapan{C}
\end{cauhoi}
\begin{cauhoi}{Câu 35}{8H2-C7-B28-D02-TN-TH-035}
Tìm toạ độ giao điểm của hai đường thẳng $(d_1): y=-x+7$ và $(d_2): y=4x+12$.
\begin{tasks}(2)
\task $(8;-1)$
\task $(-1;-8)$
\task $(1;8)$
\task $(-1;8)$
\end{tasks}
\dapan{D}
\end{cauhoi}
\begin{cauhoi}{Câu 36}{8H2-C7-B28-D02-TN-TH-036}
Tìm toạ độ giao điểm của hai đường thẳng $(d_1): y=4x+6$ và $(d_2): y=-4x+6$.
\begin{tasks}(2)
\task $(1;6)$
\task $(0;-6)$
\task $(6;0)$
\task $(0;6)$
\end{tasks}
\dapan{D}
\end{cauhoi}
\begin{cauhoi}{Câu 37}{8H2-C7-B28-D02-TN-TH-037}
Tìm toạ độ giao điểm của hai đường thẳng $(d_1): y=-2x+8$ và $(d_2): y=-4x+10$.
\begin{tasks}(2)
\task $(-1;6)$
\task $(1;6)$
\task $(6;1)$
\task $(1;-6)$
\end{tasks}
\dapan{B}
\end{cauhoi}
\begin{cauhoi}{Câu 38}{8H2-C7-B28-D02-TN-TH-038}
Tìm toạ độ giao điểm của hai đường thẳng $(d_1): y=3x-4$ và $(d_2): y=-3x+2$.
\begin{tasks}(2)
\task $(1;-1)$
\task $(1;1)$
\task $(-1;1)$
\task $(-1;-1)$
\end{tasks}
\dapan{A}
\end{cauhoi}
\begin{cauhoi}{Câu 39}{8H2-C7-B28-D02-TN-TH-039}
Tìm toạ độ giao điểm của hai đường thẳng $(d_1): y=4x+24$ và $(d_2): y=2x+14$.
\begin{tasks}(2)
\task $(4;-5)$
\task $(5;4)$
\task $(-5;4)$
\task $(-5;-4)$
\end{tasks}
\dapan{C}
\end{cauhoi}
\begin{cauhoi}{Câu 40}{8H2-C7-B28-D02-TN-TH-040}
Tìm toạ độ giao điểm của hai đường thẳng $(d_1): y=-2x-3$ và $(d_2): y=x+9$.
\begin{tasks}(2)
\task $(-4;5)$
\task $(4;5)$
\task $(-4;-5)$
\task $(5;-4)$
\end{tasks}
\dapan{A}
\end{cauhoi}
\begin{cauhoi}{Câu 41}{8H2-C7-B28-D02-TN-TH-041}
Tìm toạ độ giao điểm của hai đường thẳng $(d_1): y=-x+4$ và $(d_2): y=2x+4$.
\begin{tasks}(2)
\task $(4;0)$
\task $(1;4)$
\task $(0;4)$
\task $(0;-4)$
\end{tasks}
\dapan{C}
\end{cauhoi}
\begin{cauhoi}{Câu 42}{8H2-C7-B28-D02-TN-TH-042}
Tìm toạ độ giao điểm của hai đường thẳng $(d_1): y=-x+7$ và $(d_2): y=2x+10$.
\begin{tasks}(2)
\task $(-1;-8)$
\task $(-1;8)$
\task $(8;-1)$
\task $(1;8)$
\end{tasks}
\dapan{B}
\end{cauhoi}
\begin{cauhoi}{Câu 43}{8H2-C7-B28-D02-TN-TH-043}
Tìm toạ độ giao điểm của hai đường thẳng $(d_1): y=4x+8$ và $(d_2): y=-x-2$.
\begin{tasks}(2)
\task $(-1;0)$
\task $(-2;0)$
\task $(2;0)$
\task $(0;-2)$
\end{tasks}
\dapan{B}
\end{cauhoi}
\begin{cauhoi}{Câu 44}{8H2-C7-B28-D02-TN-TH-044}
Tìm toạ độ giao điểm của hai đường thẳng $(d_1): y=-x-4$ và $(d_2): y=2x+2$.
\begin{tasks}(2)
\task $(-1;-2)$
\task $(-2;2)$
\task $(2;-2)$
\task $(-2;-2)$
\end{tasks}
\dapan{D}
\end{cauhoi}
\begin{cauhoi}{Câu 45}{8H2-C7-B28-D02-TN-TH-045}
Tìm toạ độ giao điểm của hai đường thẳng $(d_1): y=4x+20$ và $(d_2): y=-4x-20$.
\begin{tasks}(2)
\task $(5;0)$
\task $(-4;0)$
\task $(0;-5)$
\task $(-5;0)$
\end{tasks}
\dapan{D}
\end{cauhoi}
\begin{cauhoi}{Câu 46}{8H2-C7-B28-D02-TN-TH-046}
Tìm toạ độ giao điểm của hai đường thẳng $(d_1): y=x-2$ và $(d_2): y=-3x+6$.
\begin{tasks}(2)
\task $(2;0)$
\task $(0;2)$
\task $(3;0)$
\task $(-2;0)$
\end{tasks}
\dapan{A}
\end{cauhoi}
\begin{cauhoi}{Câu 47}{8H2-C7-B28-D02-TN-TH-047}
Tìm toạ độ giao điểm của hai đường thẳng $(d_1): y=-2x+12$ và $(d_2): y=-3x+16$.
\begin{tasks}(2)
\task $(-4;4)$
\task $(5;4)$
\task $(4;-4)$
\task $(4;4)$
\end{tasks}
\dapan{D}
\end{cauhoi}
\begin{cauhoi}{Câu 48}{8H2-C7-B28-D02-TN-TH-048}
Tìm toạ độ giao điểm của hai đường thẳng $(d_1): y=3x+14$ và $(d_2): y=2x+10$.
\begin{tasks}(2)
\task $(-4;2)$
\task $(-4;-2)$
\task $(4;2)$
\task $(2;-4)$
\end{tasks}
\dapan{A}
\end{cauhoi}
\begin{cauhoi}{Câu 49}{8H2-C7-B28-D02-TN-TH-049}
Tìm toạ độ giao điểm của hai đường thẳng $(d_1): y=x+8$ và $(d_2): y=2x+8$.
\begin{tasks}(2)
\task $(0;-8)$
\task $(1;8)$
\task $(8;0)$
\task $(0;8)$
\end{tasks}
\dapan{D}
\end{cauhoi}
\begin{cauhoi}{Câu 50}{8H2-C7-B28-D02-TN-TH-050}
Tìm toạ độ giao điểm của hai đường thẳng $(d_1): y=3x-13$ và $(d_2): y=x-5$.
\begin{tasks}(2)
\task $(-1;4)$
\task $(-4;-1)$
\task $(4;1)$
\task $(4;-1)$
\end{tasks}
\dapan{D}
\end{cauhoi}
\clearpage\section*{Bảng đáp án trắc nghiệm}
\begin{center}\begin{tabular}{|c|c|c|c|c|c|c|c|c|c|}
\hline
1. C & 2. B & 3. A & 4. A & 5. B & 6. B & 7. D & 8. B & 9. C & 10. B\\\hline
11. C & 12. A & 13. B & 14. A & 15. A & 16. C & 17. B & 18. C & 19. A & 20. D\\\hline
21. C & 22. C & 23. C & 24. D & 25. D & 26. A & 27. B & 28. C & 29. B & 30. A\\\hline
31. A & 32. A & 33. C & 34. C & 35. D & 36. D & 37. B & 38. A & 39. C & 40. A\\\hline
41. C & 42. B & 43. B & 44. D & 45. D & 46. A & 47. D & 48. A & 49. D & 50. D\\\hline
\end{tabular}\end{center}
\end{document}
